\section{Pregunta 3. Firma digital con curva elíptica}
\label{sec:pregunta3}

\begin{enunciado}{(2,5 puntos)}
Utilizando Openssl cree un par de clave pública/privada con curva elíptica, elija cualquier curva, y firme un documento verificando posteriormente su firma
\end{enunciado}

\subsection{Análisis}

La criptografía asimétrica utiliza dos claves relacionadas: una privada, que solo conserva su propietario, y otra pública, que puede compartirse.

En la firma digital, el emisor firma con su clave privada y el receptor verifica con la pública correspondiente. Esto permite comprobar la integridad del documento y su vinculación con esa clave; para atribuirlo a una persona, debe conocerse de forma fiable a quién pertenece la clave pública.
La firma no oculta el contenido del documento en este caso \ref{ref:p3-asimetrica}, \ref{ref:p3-firma}.

En esta pregunta, el emisor firmará un fichero (el PDF \texttt{tarea.pdf}) con su clave privada $d_A$. El receptor dispondrá del fichero en cuestión, la firma $\sigma_i$ y la clave pública $Q_A$ para comprobarla, sin recibir la clave privada:

\begin{figure}[H]
  \centering
  \resizebox{\textwidth}{!}{% Original vector diagram: document signing with the sender's asymmetric key pair.
\begingroup
\setstretch{1}
\hyphenpenalty=10000
\newcommand{\sigblock}[3]{%
 {\fontsize{13}{15}\selectfont\bfseries\boldmath #1}\\[4pt]
 {\fontsize{10}{12}\selectfont #2}\\[3pt]
 {\fontsize{8}{10}\selectfont (#3)}%
}
\begin{tikzpicture}[x=1cm,y=1cm,
 box/.style={proceso,text width=27mm,minimum height=18mm,inner sep=6pt},
 size/.style={etiqueta,fill=none,font=\sffamily\fontsize{9}{11}\selectfont}]
 \node[box,fill=white,funcionpublica] (file) at (0,0)
   {\sigblock{$M_i$}{Original file}{any file format}};
 \node[box,funcioncripto] (sign) at (4,0)
   {\sigblock{Sign}{Sender: Alice}{file + private key}};
 \node[box,fill=white,funcionpublica] (signed) at (8,0)
   {\sigblock{$M_i,\ \sigma_i$}{File and signature}{file remains unchanged}};
 \node[box,funcioncripto] (verify) at (12,0)
   {\sigblock{Verify}{Receiver: Bob}{file; signature; public key}};
 \node[box,fill=white,funcionsecreta] (private) at (4,3)
   {\sigblock{$d_A$}{Private key}{kept by Alice}};
 \node[box,fill=white,funcionpublica] (public) at (12,3)
   {\sigblock{$Q_A$}{Public key}{shared with Bob}};
 \node[box,fill=white,funcionpublica] (result) at (12,-2.8)
   {\sigblock{Valid / Invalid}{Verification result}{accept or reject}};
 \draw[flujo] (file.east) -- (sign.west);
 \draw[flujo] (sign.east) -- (signed.west);
 \draw[flujo] (signed.east) -- (verify.west);
 \draw[flujo] (private.south) -- (sign.north);
 \draw[flujo] (public.south) -- (verify.north);
 \draw[flujo,dashed] (private.east) -- node[size,above] {Public key derivation} (public.west);
 \draw[flujo] (verify.south) -- (result.north);
 \leyendaFunciones{(5.7,-4.5)}
\end{tikzpicture}
\endgroup
}
  \caption{Firma y verificación de un fichero con par de claves de curva elíptica de un emisor. Elaboración propia a partir del tema 1.5 (\ref{ref:p3-firma}).}
  \label{fig:p3-claves-asimetricas}
\end{figure}

\clearpage
\textbf{Firma con curva elíptica:} ECC (Elliptic Curve Cryptography) opera con puntos de una curva sobre un cuerpo finito. La clave privada $d_A$ es un número secreto y la pública se obtiene como $Q_A=d_AG$, donde $G$ es un punto base público. Calcular esa multiplicación de puntos es eficiente; recuperar $d_A$ desde $G$ y $Q_A$ resulta computacionalmente inviable con los parámetros adecuados \ref{ref:p3-curvas}.

Se elegirá una curva estandarizada sin debilidades conocidas, definida sobre un cuerpo de tamaño suficiente. El punto $G$ debe generar un subgrupo de orden primo grande $n$ (el número de puntos del subgrupo), y $d_A$ elegirse de forma aleatoria y uniforme entre $1$ y $n-1$. En P-256, tanto el primo $p$ que define el cuerpo como el orden $n$ tienen 256 bits \ref{ref:p3-curvas}, \ref{ref:p3-parametros}.

Para identificar estos parámetros con números pequeños, consideremos la ecuación $y^2\equiv x^3+ax+b\pmod p$ con $p=17$, $a=1$, $b=1$ y $G=(6,6)$. En esta curva, $G$ genera un subgrupo de orden $n=18$. Si elegimos $d_A=2$, la clave pública será $Q_A=d_AG=G+G=(9,5)$, como muestra la figura. Estos valores son únicamente de ejemplo: su tamaño es pequeño y $n$ no es primo, por lo que no cumplen los requisitos de seguridad anteriores. 

\begin{figure}[H]
  \centering
  \resizebox{\textwidth}{!}{% Teaching example from topic 1.3: y^2 = x^3 + x + 1 modulo 17.
\begingroup
\setstretch{1}
\hyphenpenalty=10000
\newcommand{\ecblock}[3]{%
 {\fontsize{13}{15}\selectfont\bfseries\boldmath #1}\\[4pt]
 {\fontsize{10}{12}\selectfont #2}\\[3pt]
 {\fontsize{8}{10}\selectfont (#3)}%
}
\begin{tikzpicture}[x=1cm,y=1cm,
 box/.style={proceso,text width=30mm,minimum height=18mm,inner sep=6pt},
 size/.style={etiqueta,fill=none,font=\sffamily\fontsize{9}{11}\selectfont}]
 \node[font=\sffamily\fontsize{10}{12}\selectfont,align=center] at (1.85,5.2)
   {Small-field example\\$y^2\equiv x^3+x+1\pmod{17}$};
 \begin{scope}[x=0.22cm,y=0.22cm]
  \draw[step=1,draw=gray!20,line width=0.2pt] (0,0) grid (16,16);
  \draw[-{Latex[length=1.5mm]},draw=Grafito] (0,0)--(17.5,0) node[right,font=\scriptsize] {$x$};
  \draw[-{Latex[length=1.5mm]},draw=Grafito] (0,0)--(0,17.5) node[above,font=\scriptsize] {$y$};
  \foreach \t in {0,8,16}{
   \node[below,font=\sffamily\scriptsize] at (\t,0) {\t};
  }
  \foreach \t in {8,16}{\node[left,font=\sffamily\scriptsize] at (0,\t) {\t};}
  \foreach \x/\y in {0/1,0/16,4/1,4/16,6/6,6/11,9/5,9/12,10/5,10/12,11/0,13/1,13/16,15/5,15/12,16/4,16/13}{\fill[Grafito!60] (\x,\y) circle (1.3pt);}
  \fill[BordePublico] (6,6) circle (2.6pt);
  \node[above left,font=\sffamily\small,text=BordePublico] at (6,6) {$G$};
  \fill[BordePublico] (9,5) circle (2.6pt);
  \node[below right,font=\sffamily\small,text=BordePublico] at (9,5) {$Q_A$};
 \end{scope}
 \node[box,fill=white,funcionsecreta] (private) at (7.7,4.9)
   {\ecblock{$d_A=2$}{Private key}{secret scalar}};
 \node[box,funcioncripto] (multiply) at (7.7,1.8)
   {\ecblock{$Q_A=d_AG$}{Point multiplication}{$G+G$ in this example}};
 \node[box,fill=white,funcionpublica] (public) at (12.4,1.8)
   {\ecblock{$Q_A=(9,5)$}{Public key}{point on the same curve}};
 \coordinate (generator) at (1.32,1.32);
 \draw[flujo] (private.south) -- (multiply.north);
 \draw[flujo] (generator) -- (4.5,1.32) -- (4.5,1.8)
   -- node[size,above] {$G=(6,6)$} (multiply.west);
 \draw[flujo] (multiply.east) -- (public.west);
 \leyendaFunciones{(6.2,-1.3)}
\end{tikzpicture}
\endgroup
}
  \caption{Obtención de la clave pública mediante operaciones sobre puntos de una curva. Ejemplo del tema 1.3 (\ref{ref:p3-curvas}).}
  \label{fig:p3-curva-eliptica}
\end{figure}

\textbf{Composición y validación de la firma:} ECDSA (Elliptic Curve Digital Signature Algorithm) genera una firma $\sigma_i=(r,s)$: dos valores enteros entre $1$ y $n-1$, que vinculan el contenido del fichero con la clave privada del firmante. Para obtenerlos intervienen el resumen del fichero, la clave privada $d_A$, los parámetros públicos de la curva y un valor secreto por firma $k$. El valor $r$ procede del punto $kG$, mientras que $s$ combina el resumen, $d_A$, $r$ y $k$ mediante operaciones módulo $n$. La firma se conserva aparte del fichero y no contiene la clave privada \ref{ref:p3-ecdsa}.

En el esquema, $E$ identifica la curva y $z$ representa el resumen SHA-256 interpretado como entero. El símbolo $k^{-1}$ indica el inverso de $k$ módulo $n$. OpenSSL realiza estas operaciones internamente.

\begin{figure}[H]
  \centering
  \resizebox{\textwidth}{!}{% ECDSA signing, adapted from topic 1.5, slides 19-21.
% SHA-256 with P-256 is the choice made for this practice.
\begingroup
\setstretch{1}
\hyphenpenalty=10000
\newcommand{\ecdsablock}[3]{%
 {\fontsize{13}{15}\selectfont\bfseries\boldmath #1}\\[4pt]
 {\fontsize{10}{12}\selectfont #2}\\[3pt]
 {\fontsize{8}{10}\selectfont (#3)}%
}
\begin{tikzpicture}[x=1cm,y=1cm,
 box/.style={proceso,text width=29mm,minimum height=20mm,inner sep=6pt},
 size/.style={etiqueta,fill=none,font=\sffamily\fontsize{9}{11}\selectfont}]
 \node[box,fill=white,funcionpublica] (file) at (0,0)
   {\ecdsablock{$M_i$}{Original file}{unchanged by signing}};
 \node[box,funcioncripto] (hash) at (4.5,0)
   {\ecdsablock{SHA-256}{Secure Hash Algorithm}{hash the file bytes}};
 \node[box,fill=white,funcionpublica] (digest) at (9,0)
   {\ecdsablock{$z$}{Digest as an integer}{32 bytes; 256 bits}};

 \node[box,fill=white,funcionsecreta] (nonce) at (0,-3.4)
   {\ecdsablock{$k$}{Per-signature secret}{random; $1\leq k<n$}};
 \node[box,fill=white,funcionpublica] (params) at (4.5,-3.4)
   {\ecdsablock{$E,\ G,\ n$}{Curve parameters}{P-256; base point; order}};
 \node[box,fill=white,funcionsecreta] (private) at (9,-3.4)
   {\ecdsablock{$d_A$}{Private key}{$1\leq d_A<n$}};

 \node[box,funcioncripto] (point) at (0,-6.8)
   {\ecdsablock{$R=kG$}{Point multiplication}{$R=(x_R,y_R)$}};
 \node[box,funcioncripto] (r) at (4.5,-6.8)
   {\ecdsablock{$r$}{First signature value}{$r=x_R\bmod n$;\\retry with a new $k$ if $r=0$}};
 \node[box,funcioncripto,text width=37mm] (s) at (9,-6.8)
   {\ecdsablock{$s$}{Second signature value}{$s=k^{-1}(z+r d_A)\bmod n$;\\retry with a new $k$ if $s=0$}};
 \node[box,fill=white,funcionpublica,text width=37mm] (signature) at (9,-11)
   {\ecdsablock{$\sigma_i=(r,s)$}{Detached signature}{stored separately from $M_i$;\\neither $d_A$ nor $k$ is included}};

 \draw[flujo] (file.east) -- (hash.west);
 \draw[flujo] (hash.east) -- (digest.west);
 \draw[flujo] (params.west) -- node[size,above] {$n$} (nonce.east);
 \draw[flujo] (nonce.south) -- node[size,left] {$k$} (point.north);
 \draw[flujo] ([xshift=-8mm]params.south) -- (3.7,-5.1) -- (0.8,-5.1)
   -- node[size,right] {$E,G$} ([xshift=8mm]point.north);
 \draw[flujo] (params.south) -- node[size,right] {$n$} (r.north);
 \draw[flujo] (point.east) -- node[size,above] {$x_R$} (r.west);
 \draw[flujo] (r.east) -- node[size,above] {$r,n$} (s.west);
 \draw[flujo] (private.south) -- node[size,right] {$d_A$} (s.north);
 \draw[flujo] (digest.east) -- (12,0) -- (12,-6.8)
   -- node[size,above] {$z$} (s.east);
 \draw[flujo] (nonce.west) -- (-2.5,-3.4) -- (-2.5,-9)
   -- node[size,above,pos=0.7] {$k$} (8,-9) -- ([xshift=-10mm]s.south);
 \draw[flujo] ([xshift=10mm]s.south) -- (10,-9.5) -- (9,-9.5)
   -- node[size,right] {$r,s$} (signature.north);
 \leyendaFunciones{(4.8,-13.1)}
\end{tikzpicture}
\endgroup
}
  \caption{Generación de la firma ECDSA de un fichero. Adaptación del tema 1.5, diap.~21, con P-256 y SHA-256 (\ref{ref:p3-ecdsa}).}
  \label{fig:p3-ecdsa-firma}
\end{figure}

Para que este mecanismo sea seguro, $d_A$ y $k$ deben permanecer secretos. En la generación aleatoria, $k$ debe ser impredecible, estar entre $1$ y $n-1$ y no reutilizarse en documentos distintos, pues podría revelarse $d_A$. El verificador necesita el fichero, $(r,s)$ y una clave pública $Q_A$ de origen fiable, además de utilizar la misma curva y función de resumen que el firmante \ref{ref:p3-firma}.

La verificación comprueba el rango de $r$ y $s$, calcula el resumen del fichero recibido y verifica su relación matemática con la firma y $Q_A$. Si se modifica el contenido firmado, la comprobación debe fallar; disponer solo de la clave pública no permite generar una firma válida para el fichero modificado \ref{ref:p3-ecdsa}.

En la práctica se empleará P-256, por ser una curva estandarizada y compatible con ECDSA en OpenSSL, donde se denomina \texttt{prime256v1}. Sus parámetros públicos $(p,a,b,G,n,h)$ están fijados por el estándar: $p$ define el cuerpo, $a$ y $b$ la ecuación, $G$ es el punto base, $n$ su orden primo y $h$ el cofactor, es decir, el cociente entre el número total de puntos de la curva y $n$.
En P-256, $p$ y $n$ tienen 256 bits, $a=-3\pmod p$ y $h=1$; los valores de $b$ y las coordenadas de $G$ son constantes estandarizadas.
OpenSSL carga estos parámetros al seleccionar la curva, sin elegir un $G$ nuevo para cada clave \ref{ref:p3-curvas}, \ref{ref:p3-parametros}, \ref{ref:p3-openssl}.

\begin{table}[H]
  \centering
  \caption{Parámetros públicos fijos de P-256 (\ref{ref:p3-parametros}).}
  \label{tab:p3-parametros-p256}
  \begingroup
  \small
  \renewcommand{\arraystretch}{1.35}
  \begin{tabularx}{\textwidth}{@{}l l X@{}}
    \toprule
    \textbf{Parámetro} & \textbf{Significado} & \textbf{Valor} \\
    \midrule
$p$ & Primo del cuerpo & $2^{256}-2^{224}+2^{192}+2^{96}-1$ \\
    $a$ & Coeficiente & $-3\pmod p$ \\
    $b$ & Coeficiente & \shortstack[l]{\ttfamily\footnotesize 5ac635d8 aa3a93e7 b3ebbd55 769886bc\\
\ttfamily\footnotesize 651d06b0 cc53b0f6 3bce3c3e 27d2604b} \\
    $G_x$ & 1º Coordenada de $G$ & \shortstack[l]{\ttfamily\footnotesize 6b17d1f2 e12c4247 f8bce6e5 63a440f2\\
\ttfamily\footnotesize 77037d81 2deb33a0 f4a13945 d898c296} \\
    $G_y$ & 2º Coordenada de $G$ & \shortstack[l]{\ttfamily\footnotesize 4fe342e2 fe1a7f9b 8ee7eb4a 7c0f9e16\\
\ttfamily\footnotesize 2bce3357 6b315ece cbb64068 37bf51f5} \\
    $n$ & Orden de $G$ & \shortstack[l]{\ttfamily\footnotesize ffffffff 00000000 ffffffff ffffffff\\
\ttfamily\footnotesize bce6faad a7179e84 f3b9cac2 fc632551} \\
    $h$ & Cofactor & $1$ \\
    \bottomrule
  \end{tabularx}
  \endgroup

  \vspace{3pt}
  \parbox{\textwidth}{\footnotesize Los valores de $b$, $G_x$, $G_y$ y $n$ están en hexadecimal; las dos líneas de cada celda forman un único número. El punto base es $G=(G_x,G_y)$.}
\end{table}

Se utilizará SHA-256 como función de resumen. OpenSSL gestionará la generación de las claves y del valor $k$, así como la firma y su verificación.


\subsection{Implementación y validación}

La solución se implementó mediante tres scripts compatibles con \texttt{/bin/sh}, ejecutados desde \texttt{scripts/pregunta-3}. Se utilizó el PDF original \texttt{tarea.pdf}, de 1\,560\,658 bytes, conservado en \texttt{originals} del workspace común como fichero a firmar.
El par de claves de firma se generó de forma independiente de las claves simétricas anteriores: la privada se guardó en \texttt{secrets/ec-keys}, la pública en \texttt{metadata/public-keys} y la firma separada en \texttt{signatures}.

\textbf{Generación del par de claves:} El script \texttt{gen-keys.sh} establece los destinos y comprueba que no existan claves previas, evitando sustituir el par utilizado por firmas anteriores. Con \texttt{openssl genpkey}, selecciona \texttt{prime256v1} y genera la privada; después, \texttt{openssl pkey -pubout} obtiene la pública correspondiente \ref{ref:p3-openssl}.

Las claves se guardan en formato PEM (Privacy-Enhanced Mail), con permisos restringidos al propietario mediante \texttt{umask 077} y \texttt{chmod 600}. El script valida ambas claves con \texttt{-check} y \texttt{-pubcheck}; si falla la operación, las órdenes \texttt{trap} eliminan el par incompleto. La información mostrada al finalizar corresponde exclusivamente a la clave pública.

\begingroup
\renewcommand{\lstlistingname}{Script}
\setstretch{1}
\lstinputlisting[style=programacompleto,language=bash,
 caption={Generación y validación del par de claves de curva elíptica de P-256: gen-keys.sh.},
 label={lst:p3-gen-completo}]{recursos/listados/pregunta-3/gen-keys.sh}
\endgroup

Tras la generación de dicho par de claves, se confirmó la validez de ambas y la selección de P-256:

\noindent\begin{minipage}{\linewidth}
\begingroup
\renewcommand{\lstlistingname}{Salida y verificación}
\setstretch{1}
\lstinputlisting[style=terminal,basicstyle=\ttfamily\footnotesize,
 caption={Generación de par de claves satisfactoria y consulta de la clave pública.},
 label={lst:p3-claves-real}]{recursos/salidas/pregunta-3/generacion-claves.txt}
\endgroup
\end{minipage}

El punto público $Q_A$ (clave pública) se muestra en hexadecimal: el prefijo \texttt{04} indica la representación no comprimida, seguido de las coordenadas $x$ e $y$, de 32 bytes cada una. El campo de 256 bits no equivale a 256 bits de seguridad; la salida indica un nivel de seguridad de 128 bits.

\textbf{Firma del fichero:} El script \texttt{sign.sh} utiliza por defecto el PDF del enunciado original como fichero a firmar y la clave privada generada, aunque admite otras rutas como argumentos. Comprueba las entradas y que la clave corresponda a P-256. También evita que el destino de la firma coincida con el original o con la clave.

La orden \texttt{openssl dgst -sha256 -sign} calcula el resumen de los bytes del fichero y genera la firma ECDSA. No se calcula un resumen previo por separado. El par $(r,s)$ se guarda con codificación DER en \texttt{tarea.pdf.sig}, conservando el fichero original.
La firma se escribe directamente en su destino, sobrescribiendo una anterior; si la operación falla, se elimina la salida incompleta \ref{ref:p3-openssl}.

\begingroup
\renewcommand{\lstlistingname}{Script}
\setstretch{1}
\lstinputlisting[style=programacompleto,language=bash,
 caption={Firma del fichero con ECDSA y SHA-256: sign.sh.},
 label={lst:p3-sign-completo}]{recursos/listados/pregunta-3/sign.sh}
\endgroup

La ejecución produjo una firma de 70 bytes para el fichero de 1\,560\,658 bytes:

\noindent\begin{minipage}{\linewidth}
\begingroup
\renewcommand{\lstlistingname}{Salida y verificación}
\setstretch{1}
\lstinputlisting[style=terminal,basicstyle=\ttfamily\footnotesize,
 caption={Ejecución de la firma y tamaños de los archivos.},
 label={lst:p3-firma-real}]{recursos/salidas/pregunta-3/firma.txt}
\endgroup
\end{minipage}

A continuación se inspeccionó su estructura mediante \texttt{openssl asn1parse}, que interpreta la codificación ASN.1 (Abstract Syntax Notation One):

\noindent\begin{minipage}{\linewidth}
\begingroup
\renewcommand{\lstlistingname}{Salida y verificación}
\setstretch{1}
\lstinputlisting[style=terminal,basicstyle=\ttfamily\footnotesize,
 caption={Consulta de los datos que componen la firma.},
 label={lst:p3-componentes-real}]{recursos/salidas/pregunta-3/componentes-firma.txt}
\endgroup
\end{minipage}

Los dos campos \texttt{INTEGER} corresponden, en ese orden, a $r$ y $s$. En esta ejecución, cada entero ocupa 32 bytes; los 6 bytes restantes pertenecen a la estructura DER, hasta completar los 70 bytes del archivo. El tamaño codificado puede variar entre firmas y no representa el tamaño del fichero firmado.

\textbf{Verificación y prueba de modificación:} El script \texttt{verify.sh} recibe el fichero original, la clave pública y la firma; no necesita la privada.
Comprueba que los tres archivos existan y que la pública sea P-256. Después utiliza \texttt{openssl dgst -sha256 -verify} para recalcular el resumen y verificar la firma.
Si OpenSSL termina correctamente, muestra \texttt{ACCEPT}; en caso contrario, muestra \texttt{REJECT} y devuelve un código de error, sin eliminar ninguna entrada \ref{ref:p3-openssl}.

\begingroup
\renewcommand{\lstlistingname}{Script}
\setstretch{1}
\lstinputlisting[style=programacompleto,language=bash,
 caption={Verificación de la firma con la clave pública: verify.sh.},
 label={lst:p3-verify-completo}]{recursos/listados/pregunta-3/verify.sh}
\endgroup

La primera verificación de la firma finalizó correctamente:

\noindent\begin{minipage}{\linewidth}
\begingroup
\renewcommand{\lstlistingname}{Salida y verificación}
\setstretch{1}
\lstinputlisting[style=terminal,basicstyle=\ttfamily\footnotesize,
 caption={Aceptación de la firma del fichero original con la clave pública.},
 label={lst:p3-verificacion-original}]{recursos/salidas/pregunta-3/verificacion-original.txt}
\endgroup
\end{minipage}

Para comprobar la detección de cambios, se creó \texttt{tarea-modified.pdf} como copia del original y se añadió una línea de texto a sus bytes.
Se verificó esta copia con la misma clave pública y la firma del original, sin generar una firma nueva:

\noindent\begin{minipage}{\linewidth}
\begingroup
\renewcommand{\lstlistingname}{Salida y verificación}
\setstretch{1}
\lstinputlisting[style=terminal,basicstyle=\ttfamily\footnotesize,
 columns=fixed,basewidth=0.5em,breakindent=1em,breakatwhitespace=false,
 caption={Modificación de una copia y rechazo de su verificación.},
 label={lst:p3-verificacion-modificado}]{recursos/salidas/pregunta-3/verificacion-modificado.txt}
\endgroup
\end{minipage}

La salida registra \texttt{Verification failure}, los diagnósticos internos emitidos por OpenSSL y el mensaje \texttt{REJECT} del script.
Se observa el rechazo de la combinación formada por la copia modificada y la firma original, asegurando que dicho fichero no es íntegro.

Finalmente, se repitió la verificación con el fichero original y la misma clave pública y firma:

\noindent\begin{minipage}{\linewidth}
\begingroup
\renewcommand{\lstlistingname}{Salida y verificación}
\setstretch{1}
\lstinputlisting[style=terminal,basicstyle=\ttfamily\footnotesize,
 caption={Verificación correcta del fichero original después de la prueba de modificación.},
 label={lst:p3-verificacion-final}]{recursos/salidas/pregunta-3/verificacion-final.txt}
\endgroup
\end{minipage}

Este procedimiento completa la correcta generación del par de claves pública/privada de curva elíptica P-256, la firma del fichero y su verificación, incluyendo el rechazo del contenido modificado.
